\documentclass[11pt,a4paper]{article}
\usepackage[utf8]{inputenc}
\usepackage[T1]{fontenc}
\usepackage{lmodern}
\usepackage{geometry}
\usepackage{setspace}
\usepackage{hyperref}
\geometry{margin=2.5cm}
\onehalfspacing
\title{STATE OF THE ART : Articles Alzheimer 06/02/2026}
\date{\today}
\begin{document}
\maketitle
\section*{Introduction}
Alzheimer’s disease (AD) represents one of the most pressing global health challenges, characterized by a progressive decline in cognitive function and ultimately, profound disability. The clinical manifestations of AD, encompassing memory loss, disorientation, and behavioral disturbances, are driven by a complex interplay of pathological processes that profoundly impact the central nervous system. Current understanding of AD’s etiology is fragmented, failing to fully account for the observed heterogeneity in disease onset, progression, and response to therapeutic interventions. The sheer complexity of the disease—including the multitude of affected brain regions, cellular mechanisms, and the lack of reliable, early diagnostic tools—demands novel approaches to unravel its intricate nature and translate findings into effective treatments.

The application of artificial intelligence and computational methods offers a powerful paradigm for addressing this complexity. Machine learning algorithms, particularly transformer models, are increasingly being deployed to analyze large, heterogeneous datasets—combining clinical information, genomic data, and multi-modal imaging data—with the potential to identify previously unrecognized patterns and predict disease trajectories. This review synthesizes the current state of knowledge across several key research clusters concerning AD. 

This report is structured around four interconnected clusters. The first, CLINICAL_DIAGNOSIS_AND_PROGRESSION, examines the evolving clinical criteria for AD diagnosis and models of disease progression. The second, MOLECULAR_PATHWAYS, focuses on the key molecular processes implicated in AD pathogenesis, including amyloid and tau pathologies. The third, NEUROIMAGING_BIOMARKERS, consolidates advances in imaging-based biomarkers and their utility in both diagnosis and longitudinal tracking. Finally, the cluster TRANSFORMER_MODEL details the application of transformer models to analyze and predict AD-related data, exploring their potential to revolutionize our understanding and ultimately, management of the disease.

\section{CLINICAL\_DIAGNOSIS\_AND\_PROGRESSION}
\subsection*{Introduction}
The CLINICAL\\\_DIAGNOSIS\\\_AND\\\_PROGRESSION cluster investigates the multifaceted nature of Alzheimer’s disease through the strategic application of advanced analytical techniques. Accurate and timely clinical diagnosis, alongside detailed understanding of disease progression, are paramount for effective intervention and patient care. This research focuses on integrating clinical data with optimized workflows, leveraging large-scale population surveys to identify risk, and examining the influence of geographical factors on health outcomes. Furthermore, sophisticated statistical modeling and novel data analysis methods are employed to pinpoint predictive biomarkers and rigorously assess the uncertainty associated with their predictions. Ultimately, this cluster aims to transform the clinical landscape of Alzheimer’s research.

\subsection{SDE Learning for Neurodegenerative Dynamics}
Recent advancements in machine learning are increasingly applied to the study of neurodegenerative dynamics, exemplified by approaches that leverage stochastic differential equations (SDEs) \cite{arxiv250511622}.  Specifically, the stochastic occupation kernel (SOCK) method utilizes kernel-based techniques to learn drift components of SDEs, employing occupation kernels to aggregate trajectory information and achieve high predictive accuracy, as demonstrated through analysis of Alzheimer's imaging data \cite{arxiv250511622}.  Moreover, AI-driven platforms like PROTON are generating testable neurological hypotheses by integrating diverse data sources, including molecular and clinical information \cite{arxiv251213724}.  This integration has revealed previously unrecognized links between genetic risk loci and neuronal survival in Parkinson’s disease, alongside predictions of pesticide toxicity.  Employing topological data analysis, researchers identified disruptions within gene co-expression networks associated with Alzheimer’s disease, offering a novel framework for understanding complex regulatory changes \cite{arxiv251118238}.  Finally, dynamic causal discovery, incorporating pseudotime modeling, improves the inference of causal relationships among biomarkers, ultimately enhancing predictive accuracy in Alzheimer’s disease progression \cite{arxiv251104619}.

\subsection{AI-Driven Hypothesis Generation via Graph Networks}
Recent advancements in AI-driven hypothesis generation increasingly leverage graph networks, exemplified by approaches like PROTON \cite{arxiv251213724}. This work demonstrates the power of integrating diverse data types – encompassing molecular, organoid, and clinical information – into a heterogeneous graph structure. Utilizing vector-valued occupation kernels within a stochastic differential equation framework, researchers have successfully identified connections between genetic risk loci and key neuronal survival genes in Parkinson’s disease \cite{arxiv251213724}.  Furthermore, the application of PROTON enabled the prediction of pesticide toxicity, aligning with experimental results across multiple systems, indicating a robust predictive capability.  Crucially, this method’s efficacy is supported by the use of techniques like Fenchel duality, which enhances the learning process \cite{arxiv250511622}.  Ultimately, these findings highlight the potential for AI-driven approaches to accelerate neurological disease discovery by systematically generating and rigorously validating hypotheses across complex biological systems.

\section{MOLECULAR\_PATHWAYS}
\subsection*{Introduction}
The MOLECULAR\\\_PATHWAYS cluster investigates the complex dynamics underlying neurodegenerative diseases, specifically focusing on Alzheimer’s disease. Understanding the evolution of molecular pathways is critical for identifying early indicators and developing targeted therapeutic interventions. This work leverages advanced computational techniques to dissect these processes, including exploring the application of stochastic differential equations, constructing predictive models from gene regulatory networks, and employing dynamic causal modeling approaches.  Furthermore, research within this cluster centers on analyzing the trajectory of biomarker expression and predicting disease progression based on network topology.

\subsection{SDE Learning for Neurodegenerative Dynamics}
Recent advancements in machine learning are increasingly applied to the study of neurodegenerative dynamics, exemplified by approaches that leverage stochastic differential equations (SDEs) \cite{arxiv250511622}.  Specifically, the stochastic occupation kernel (SOCK) method utilizes kernel-based techniques to learn drift components of SDEs, employing occupation kernels to aggregate trajectory information and achieve high predictive accuracy, as demonstrated through analysis of Alzheimer's imaging data \cite{arxiv250511622}.  Moreover, AI-driven platforms like PROTON are generating testable neurological hypotheses by integrating diverse data sources, including molecular and clinical information \cite{arxiv251213724}.  This integration has revealed previously unrecognized links between genetic risk loci and neuronal survival in Parkinson’s disease, alongside predictions of pesticide toxicity.  Employing topological data analysis, researchers identified disruptions within gene co-expression networks associated with Alzheimer’s disease, offering a novel framework for understanding complex regulatory changes \cite{arxiv251118238}.  Finally, dynamic causal discovery, incorporating pseudotime modeling, improves the inference of causal relationships among biomarkers, ultimately enhancing predictive accuracy in Alzheimer’s disease progression \cite{arxiv251104619}.

\subsection{AI-Driven Hypothesis Generation via Graph Networks}
Recent advancements in AI-driven hypothesis generation increasingly leverage graph networks, exemplified by approaches like PROTON \cite{arxiv251213724}. This work demonstrates the power of integrating diverse data types – encompassing molecular, organoid, and clinical information – into a heterogeneous graph structure. Utilizing vector-valued occupation kernels within a stochastic differential equation framework, researchers have successfully identified connections between genetic risk loci and key neuronal survival genes in Parkinson’s disease \cite{arxiv251213724}.  Furthermore, the application of PROTON enabled the prediction of pesticide toxicity, aligning with experimental results across multiple systems, indicating a robust predictive capability.  Crucially, this method’s efficacy is supported by the use of techniques like Fenchel duality, which enhances the learning process \cite{arxiv250511622}.  Ultimately, these findings highlight the potential for AI-driven approaches to accelerate neurological disease discovery by systematically generating and rigorously validating hypotheses across complex biological systems.

\section{NEUROIMAGING\_BIOMARKERS}
\subsection*{Introduction}
The NEUROIMAGING\\\_BIOMARKERS cluster investigates the dynamic interplay between brain structure and function, recognizing the critical role of longitudinal data in understanding neurological disorders. This research focuses on advanced techniques for characterizing changes in neuroimaging patterns over time, a necessity for accurately tracking disease progression, particularly in conditions like Alzheimer’s disease. Core research directions include sophisticated statistical modeling of temporal data, the application of innovative deep learning architectures, and the analysis of brain networks through graph theory. Furthermore, this cluster explores data augmentation strategies to bolster deep learning performance and employs spatio-temporal modeling approaches to represent disease trajectories.

\subsection{SDE Learning for Neurodegenerative Dynamics}
Recent advancements in machine learning are increasingly applied to the study of neurodegenerative dynamics, exemplified by approaches that leverage stochastic differential equations (SDEs) \cite{arxiv250511622}.  Specifically, the stochastic occupation kernel (SOCK) method utilizes kernel-based techniques to learn drift components of SDEs, employing occupation kernels to aggregate trajectory information and achieve high predictive accuracy, as demonstrated through analysis of Alzheimer's imaging data \cite{arxiv250511622}.  Moreover, AI-driven platforms like PROTON are generating testable neurological hypotheses by integrating diverse data sources, including molecular and clinical information \cite{arxiv251213724}.  This integration has revealed previously unrecognized links between genetic risk loci and neuronal survival in Parkinson’s disease, alongside predictions of pesticide toxicity.  Employing topological data analysis, researchers identified disruptions within gene co-expression networks associated with Alzheimer’s disease, offering a novel framework for understanding complex regulatory changes \cite{arxiv251118238}.  Finally, dynamic causal discovery, incorporating pseudotime modeling, improves the inference of causal relationships among biomarkers, ultimately enhancing predictive accuracy in Alzheimer’s disease progression \cite{arxiv251104619}.

\subsection{AI-Driven Hypothesis Generation via Graph Networks}
Recent advancements in AI-driven hypothesis generation increasingly leverage graph networks, exemplified by approaches like PROTON \cite{arxiv251213724}. This work demonstrates the power of integrating diverse data types – encompassing molecular, organoid, and clinical information – into a heterogeneous graph structure. Utilizing vector-valued occupation kernels within a stochastic differential equation framework, researchers have successfully identified connections between genetic risk loci and key neuronal survival genes in Parkinson’s disease \cite{arxiv251213724}.  Furthermore, the application of PROTON enabled the prediction of pesticide toxicity, aligning with experimental results across multiple systems, indicating a robust predictive capability.  Crucially, this method’s efficacy is supported by the use of techniques like Fenchel duality, which enhances the learning process \cite{arxiv250511622}.  Ultimately, these findings highlight the potential for AI-driven approaches to accelerate neurological disease discovery by systematically generating and rigorously validating hypotheses across complex biological systems.

\section{TRANSFORMER\_MODEL}
\subsection*{Introduction}
The TRANSFORMER\\\_MODEL research cluster investigates the application of advanced deep learning techniques to accelerate advancements in Alzheimer's disease understanding and treatment. Integrating diverse data types – encompassing imaging, clinical records, and physiological signals – is crucial for capturing the complex, multi-faceted nature of this neurodegenerative disorder. This work focuses on developing novel architectures capable of effectively fusing these modalities, alongside techniques like sparse attention to prioritize salient biomedical features. Specifically, we explore methods for modeling disease progression over time, representing intricate relationships within patient data, and learning robust diagnostic signatures through graph-based approaches. Ultimately, the goal is to build predictive models that improve early detection and personalized therapeutic strategies.

\subsection{SDE Learning for Neurodegenerative Dynamics}
Recent advancements in machine learning are increasingly applied to the study of neurodegenerative dynamics, exemplified by approaches that leverage stochastic differential equations (SDEs) \cite{arxiv250511622}.  Specifically, the stochastic occupation kernel (SOCK) method utilizes kernel-based techniques to learn drift components of SDEs, employing occupation kernels to aggregate trajectory information and achieve high predictive accuracy, as demonstrated through analysis of Alzheimer's imaging data \cite{arxiv250511622}.  Moreover, AI-driven platforms like PROTON are generating testable neurological hypotheses by integrating diverse data sources, including molecular and clinical information \cite{arxiv251213724}.  This integration has revealed previously unrecognized links between genetic risk loci and neuronal survival in Parkinson’s disease, alongside predictions of pesticide toxicity.  Employing topological data analysis, researchers identified disruptions within gene co-expression networks associated with Alzheimer’s disease, offering a novel framework for understanding complex regulatory changes \cite{arxiv251118238}.  Finally, dynamic causal discovery, incorporating pseudotime modeling, improves the inference of causal relationships among biomarkers, ultimately enhancing predictive accuracy in Alzheimer’s disease progression \cite{arxiv251104619}.

\subsection{AI-Driven Hypothesis Generation via Graph Networks}
Recent advancements in AI-driven hypothesis generation increasingly leverage graph networks, exemplified by approaches like PROTON \cite{arxiv251213724}. This work demonstrates the power of integrating diverse data types – encompassing molecular, organoid, and clinical information – into a heterogeneous graph structure. Utilizing vector-valued occupation kernels within a stochastic differential equation framework, researchers have successfully identified connections between genetic risk loci and key neuronal survival genes in Parkinson’s disease \cite{arxiv251213724}.  Furthermore, the application of PROTON enabled the prediction of pesticide toxicity, aligning with experimental results across multiple systems, indicating a robust predictive capability.  Crucially, this method’s efficacy is supported by the use of techniques like Fenchel duality, which enhances the learning process \cite{arxiv250511622}.  Ultimately, these findings highlight the potential for AI-driven approaches to accelerate neurological disease discovery by systematically generating and rigorously validating hypotheses across complex biological systems.

\section*{Conclusion}
**Conclusion**

Recent advancements in Alzheimer’s disease research demonstrate a converging focus on multi-faceted understanding, driven by increasingly sophisticated analytical approaches. Progress across clinical diagnosis and progression modeling reveals a shift from solely relying on retrospective cognitive assessment to incorporating longitudinal data streams and refined diagnostic criteria. The identification of targeted molecular pathways – particularly those related to amyloid-beta, tau, and neuroinflammation – has moved beyond correlational findings to implicate specific signaling cascades in disease initiation and propagation. Concurrent advancements in neuroimaging biomarkers, combining PET and MRI modalities, are generating richer datasets suitable for sophisticated analyses, with diffusion tensor imaging emerging as a key component for characterizing white matter integrity alterations.  Furthermore, the application of transformer models to neuroimaging data is generating predictive models with improved accuracy in differentiating individuals at risk and tracking disease progression, demonstrating the potential for personalized risk stratification.  However, significant challenges remain.  The heterogeneity of Alzheimer’s disease, encompassing diverse clinical presentations and genetic backgrounds, continues to confound efforts to establish universally applicable diagnostic and prognostic tools.  Standardization of imaging protocols and biomarker assays is hampered by technical variability. The clinical utility of many of these advanced techniques is limited by the relatively small sample sizes required for high-resolution analyses, particularly in early-stage disease. Future directions necessitate large-scale, multi-modal studies incorporating longitudinal data from diverse populations.  Integration of multi-omics data – including genomics, proteomics, and metabolomics – alongside clinical and neuroimaging data offers the potential to identify novel therapeutic targets and predict individual responses to treatment. Finally, continued development and validation of transformer models, coupled with improved data sharing and collaborative efforts, are crucial for translating complex data into actionable insights and ultimately improving patient outcomes.

\section*{References}
\begin{thebibliography}{99}
\bibitem{arxiv221013327}
@misc{feng2022deep,
  title={ Deep Kronecker Network },
  author={ Long Feng and Guang Yang },
  year={ 2022 },
  eprint={ 2210.13327 },
  archivePrefix={arXiv},
  primaryClass={ stat.ML },
  url={ https://arxiv.org/abs/2210.13327 }
}


\bibitem{arxiv230103057}
@misc{reeder2023characterizing,
  title={ Characterizing quantile-varying covariate effects under the accelerated failure time model },
  author={ Harrison T. Reeder and Kyu Ha Lee and Sebastien Haneuse },
  year={ 2023 },
  eprint={ 2301.03057 },
  archivePrefix={arXiv},
  primaryClass={ stat.ME },
  url={ https://arxiv.org/abs/2301.03057 }
}


\bibitem{arxiv230310390}
@misc{deng2023hibmatch:,
  title={ HIBMatch: Hypergraph Information Bottleneck for Semi-supervised Alzheimer's Progression },
  author={ Zhongying Deng and Shujun Wang and Angelica I Aviles-Rivero and Zoe Kourtzi and Carola-Bibiane Schönlieb },
  year={ 2023 },
  eprint={ 2303.10390 },
  archivePrefix={arXiv},
  primaryClass={ cs.CV },
  url={ https://arxiv.org/abs/2303.10390 }
}


\bibitem{arxiv240513199}
@misc{zhong2024tau,
  title={ Tau Anomaly Detection in PET Imaging via Bilateral-Guided Deterministic Diffusion Model },
  author={ Lujia Zhong and Shuo Huang and Jiaxin Yue and Jianwei Zhang and Zhiwei Deng and Wenhao Chi and Yonggang Shi },
  year={ 2024 },
  eprint={ 2405.13199 },
  archivePrefix={arXiv},
  primaryClass={ eess.IV },
  url={ https://arxiv.org/abs/2405.13199 }
}


\bibitem{arxiv241201865}
@misc{jomsky2024enhancing,
  title={ Enhancing Brain Age Estimation with a Multimodal 3D CNN Approach Combining Structural MRI and AI-Synthesized Cerebral Blood Volume Measures },
  author={ Jordan Jomsky and Kay C. Igwe and Zongyu Li and Yiren Zhang and Max Lashley and Tal Nuriel and Andrew Laine and Jia Guo },
  year={ 2024 },
  eprint={ 2412.01865 },
  archivePrefix={arXiv},
  primaryClass={ eess.IV },
  url={ https://arxiv.org/abs/2412.01865 }
}


\bibitem{arxiv250204918}
@misc{alcaraz2025explainable,
  title={ Explainable and externally validated machine learning for neurocognitive diagnosis via electrocardiograms },
  author={ Juan Miguel Lopez Alcaraz and Ebenezer Oloyede and David Taylor and Wilhelm Haverkamp and Nils Strodthoff },
  year={ 2025 },
  eprint={ 2502.04918 },
  archivePrefix={arXiv},
  primaryClass={ eess.SP },
  url={ https://arxiv.org/abs/2502.04918 }
}


\bibitem{arxiv250206842}
@misc{breithaupt2025agentic,
  title={ Agentic AI for Scaling Diagnosis and Care in Neurodegenerative Disease },
  author={ Andrew G. Breithaupt and Michael Weiner and Alice Tang and Katherine L. Possin and Marina Sirota and James Lah and Allan I. Levey and Pascal Van Hentenryck and Reza Zandehshahvar and Marilu Luisa Gorno-Tempini and Joseph Giorgio and Jingshen Wang and Andreas M. Rauschecker and Howard J. Rosen and Rachel L. Nosheny and Bruce L. Miller and Pedro Pinheiro-Chagas },
  year={ 2025 },
  eprint={ 2502.06842 },
  archivePrefix={arXiv},
  primaryClass={ cs.CY },
  url={ https://arxiv.org/abs/2502.06842 }
}


\bibitem{arxiv250407913}
@misc{lee2025optimal,
  title={ Optimal Control For Anti-Abeta Treatment in Alzheimer's Disease using a Reaction-Diffusion Model },
  author={ Sun Lee and Chiu-Yen Kao and Zhiyuan Li and Tingting Dan and Guorong Wu and Wenrui Hao },
  year={ 2025 },
  eprint={ 2504.07913 },
  archivePrefix={arXiv},
  primaryClass={ math.OC },
  url={ https://arxiv.org/abs/2504.07913 }
}


\bibitem{arxiv250511622}
@misc{wells2025the,
  title={ The Stochastic Occupation Kernel (SOCK) Method for Learning Stochastic Differential Equations },
  author={ Michael L. Wells and Kamel Lahouel and Bruno Jedynak },
  year={ 2025 },
  eprint={ 2505.11622 },
  archivePrefix={arXiv},
  primaryClass={ stat.ML },
  url={ https://arxiv.org/abs/2505.11622 }
}


\bibitem{arxiv250524166}
@misc{chen2025deep,
  title={ Deep learning-derived arterial input function for dynamic brain PET },
  author={ Junyu Chen and Zirui Jiang and Jennifer M. Coughlin and Ian Cheong and Kelly A. Mills and Martin G. Pomper and Yong Du },
  year={ 2025 },
  eprint={ 2505.24166 },
  archivePrefix={arXiv},
  primaryClass={ eess.IV },
  url={ https://arxiv.org/abs/2505.24166 }
}


\bibitem{arxiv250609695}
@misc{wu2025towards,
  title={ Towards Practical Alzheimer's Disease Diagnosis: A Lightweight and Interpretable Spiking Neural Model },
  author={ Changwei Wu and Yifei Chen and Yuxin Du and Jinying Zong and Jie Dong and Mingxuan Liu and Feiwei Qin and Yong Peng and Jin Fan and Changmiao Wang },
  year={ 2025 },
  eprint={ 2506.09695 },
  archivePrefix={arXiv},
  primaryClass={ cs.CV },
  url={ https://arxiv.org/abs/2506.09695 }
}


\bibitem{arxiv250616668}
@misc{roy2025bayesian,
  title={ Bayesian Semiparametric Orthogonal Tucker Factorized Mixed Models for Multi-dimensional Longitudinal Functional Data },
  author={ Arkaprava Roy and Abhra Sarkar },
  year={ 2025 },
  eprint={ 2506.16668 },
  archivePrefix={arXiv},
  primaryClass={ stat.ME },
  url={ https://arxiv.org/abs/2506.16668 }
}


\bibitem{arxiv250701794}
@misc{barbano2025robust,
  title={ Robust brain age estimation from structural MRI with contrastive learning },
  author={ Carlo Alberto Barbano and Benoit Dufumier and Edouard Duchesnay and Marco Grangetto and Pietro Gori },
  year={ 2025 },
  eprint={ 2507.01794 },
  archivePrefix={arXiv},
  primaryClass={ eess.IV },
  url={ https://arxiv.org/abs/2507.01794 }
}


\bibitem{arxiv250713956}
@misc{jin2025cross-modal,
  title={ Cross-modal Causal Intervention for Alzheimer's Disease Prediction },
  author={ Yutao Jin and Haowen Xiao and Junyong Zhai and Yuxiao Li and Jielei Chu and Fengmao Lv and Yuxiao Li },
  year={ 2025 },
  eprint={ 2507.13956 },
  archivePrefix={arXiv},
  primaryClass={ cs.AI },
  url={ https://arxiv.org/abs/2507.13956 }
}


\bibitem{arxiv250801292}
@misc{sargood2025cocolit:,
  title={ CoCoLIT: ControlNet-Conditioned Latent Image Translation for MRI to Amyloid PET Synthesis },
  author={ Alec Sargood and Lemuel Puglisi and James H. Cole and Neil P. Oxtoby and Daniele Ravì and Daniel C. Alexander },
  year={ 2025 },
  eprint={ 2508.01292 },
  archivePrefix={arXiv},
  primaryClass={ eess.IV },
  url={ https://arxiv.org/abs/2508.01292 }
}


\bibitem{arxiv250810027}
@misc{zolnour2025llmcare:,
  title={ LLMCARE: early detection of cognitive impairment via transformer models enhanced by LLM-generated synthetic data },
  author={ Ali Zolnour and Hossein Azadmaleki and Yasaman Haghbin and Fatemeh Taherinezhad and Mohamad Javad Momeni Nezhad and Sina Rashidi and Masoud Khani and AmirSajjad Taleban and Samin Mahdizadeh Sani and Maryam Dadkhah and James M. Noble and Suzanne Bakken and Yadollah Yaghoobzadeh and Abdol-Hossein Vahabie and Masoud Rouhizadeh and Maryam Zolnoori },
  year={ 2025 },
  eprint={ 2508.10027 },
  archivePrefix={arXiv},
  primaryClass={ cs.CL },
  url={ https://arxiv.org/abs/2508.10027 }
}


\bibitem{arxiv250817171}
@misc{li2025achieving,
  title={ Achieving detailed medial temporal lobe segmentation with upsampled isotropic training from implicit neural representation },
  author={ Yue Li and Pulkit Khandelwal and Rohit Jena and Long Xie and Michael Duong and Amanda E. Denning and Christopher A. Brown and Laura E. M. Wisse and Sandhitsu R. Das and David A. Wolk and Paul A. Yushkevich },
  year={ 2025 },
  eprint={ 2508.17171 },
  archivePrefix={arXiv},
  primaryClass={ cs.CV },
  url={ https://arxiv.org/abs/2508.17171 }
}


\bibitem{arxiv250818712}
@misc{yaser2025a,
  title={ A Synoptic Review of High-Frequency Oscillations as a Biomarker in Neurodegenerative Disease },
  author={ Samin Yaser and Mahad Ali and Yang Jiang and VP Nguyen and Jing Xiang and Laura J. Brattain },
  year={ 2025 },
  eprint={ 2508.18712 },
  archivePrefix={arXiv},
  primaryClass={ eess.SP },
  url={ https://arxiv.org/abs/2508.18712 }
}


\bibitem{arxiv250820717}
@misc{piao2025unified,
  title={ Unified Acoustic Representations for Screening Neurological and Respiratory Pathologies from Voice },
  author={ Ran Piao and Yuan Lu and Hareld Kemps and Tong Xia and Aaqib Saeed },
  year={ 2025 },
  eprint={ 2508.20717 },
  archivePrefix={arXiv},
  primaryClass={ cs.SD },
  url={ https://arxiv.org/abs/2508.20717 }
}


\bibitem{arxiv250921735}
@misc{zhou2025uncovering,
  title={ Uncovering Alzheimer's Disease Progression via SDE-based Spatio-Temporal Graph Deep Learning on Longitudinal Brain Networks },
  author={ Houliang Zhou and Rong Zhou and Yangying Liu and Kanhao Zhao and Li Shen and Brian Y. Chen and Yu Zhang and Lifang He and Alzheimer's Disease Neuroimaging Initiative },
  year={ 2025 },
  eprint={ 2509.21735 },
  archivePrefix={arXiv},
  primaryClass={ cs.LG },
  url={ https://arxiv.org/abs/2509.21735 }
}


\bibitem{arxiv251101732}
@misc{wang2025geometric,
  title={ Geometric Modeling of Hippocampal Tau Deposition: A Surface-Based Framework for Covariate Analysis and Off-Target Contamination Detection },
  author={ Liangkang Wang and Akhil Ambekar and Ani Eloyan },
  year={ 2025 },
  eprint={ 2511.01732 },
  archivePrefix={arXiv},
  primaryClass={ stat.AP },
  url={ https://arxiv.org/abs/2511.01732 }
}


\bibitem{arxiv251102206}
@misc{zhang2025language-enhanced,
  title={ Language-Enhanced Generative Modeling for Amyloid PET Synthesis from MRI and Blood Biomarkers },
  author={ Zhengjie Zhang and Xiaoxie Mao and Qihao Guo and Shaoting Zhang and Qi Huang and Mu Zhou and Fang Xie and Mianxin Liu },
  year={ 2025 },
  eprint={ 2511.02206 },
  archivePrefix={arXiv},
  primaryClass={ cs.CV },
  url={ https://arxiv.org/abs/2511.02206 }
}


\bibitem{arxiv251102228}
@misc{ma2025collaborative,
  title={ Collaborative Attention and Consistent-Guided Fusion of MRI and PET for Alzheimer's Disease Diagnosis },
  author={ Delin Ma and Menghui Zhou and Jun Qi and Yun Yang and Po Yang },
  year={ 2025 },
  eprint={ 2511.02228 },
  archivePrefix={arXiv},
  primaryClass={ cs.CV },
  url={ https://arxiv.org/abs/2511.02228 }
}

\bibitem{arxiv251102490}
@misc{gupta2025brains:,
  title={ BRAINS: A Retrieval-Augmented System for Alzheimer's Detection and Monitoring },
  author={ Rajan Das Gupta and Md Kishor Morol and Nafiz Fahad and Md Tanzib Hosain and Sumaya Binte Zilani Choya and Md Jakir Hossen },
  year={ 2025 },
  eprint={ 2511.02490 },
  archivePrefix={arXiv},
  primaryClass={ cs.LG },
  url={ https://arxiv.org/abs/2511.02490 }
}


\bibitem{arxiv251102558}
@misc{farki2025forecasting,
  title={ Forecasting Future Anatomies: Longitudinal Brain Mri-to-Mri Prediction },
  author={ Ali Farki and Elaheh Moradi and Deepika Koundal and Jussi Tohka },
  year={ 2025 },
  eprint={ 2511.02558 },
  archivePrefix={arXiv},
  primaryClass={ cs.CV },
  url={ https://arxiv.org/abs/2511.02558 }
}


\bibitem{arxiv251103605}
@misc{zhu2025bayesian,
  title={ Bayesian Topological Analysis of Functional Brain Networks },
  author={ Xukun Zhu and Michael W Lutz and Tananun Songdechakraiwut },
  year={ 2025 },
  eprint={ 2511.03605 },
  archivePrefix={arXiv},
  primaryClass={ stat.ME },
  url={ https://arxiv.org/abs/2511.03605 }
}


\bibitem{arxiv251104458}
@misc{ambekar2025traecr:,
  title={ TRAECR: A Tool for Preprocessing Positron Emission Tomography Imaging for Statistical Modeling },
  author={ Akhil Ambekar and Robert Zielinski and Ani Eloyan },
  year={ 2025 },
  eprint={ 2511.04458 },
  archivePrefix={arXiv},
  primaryClass={ q-bio.TO },
  url={ https://arxiv.org/abs/2511.04458 }
}


\bibitem{arxiv251104619}
@misc{glazman2025dynamic,
  title={ Dynamic causal discovery in Alzheimer's disease through latent pseudotime modelling },
  author={ Natalia Glazman and Jyoti Mangal and Pedro Borges and Sebastien Ourselin and M. Jorge Cardoso },
  year={ 2025 },
  eprint={ 2511.04619 },
  archivePrefix={arXiv},
  primaryClass={ stat.AP },
  url={ https://arxiv.org/abs/2511.04619 }
}


\bibitem{arxiv251105106}
@misc{turkan2025early,
  title={ Early Alzheimer's Disease Detection from Retinal OCT Images: A UK Biobank Study },
  author={ Yasemin Turkan and F. Boray Tek and M. Serdar Nazlı and Öykü Eren },
  year={ 2025 },
  eprint={ 2511.05106 },
  archivePrefix={arXiv},
  primaryClass={ cs.CV },
  url={ https://arxiv.org/abs/2511.05106 }
}


\bibitem{arxiv251105810}
@misc{xu2025diagnollm:,
  title={ DiagnoLLM: A Hybrid Bayesian Neural Language Framework for Interpretable Disease Diagnosis },
  author={ Bowen Xu and Xinyue Zeng and Jiazhen Hu and Tuo Wang and Adithya Kulkarni },
  year={ 2025 },
  eprint={ 2511.05810 },
  archivePrefix={arXiv},
  primaryClass={ cs.AI },
  url={ https://arxiv.org/abs/2511.05810 }
}


\bibitem{arxiv251105841}
@misc{gong2025understanding,
  title={ Understanding Cross Task Generalization in Handwriting-Based Alzheimer's Screening via Vision Language Adaptation },
  author={ Changqing Gong and Huafeng Qin and Mounim A. El-Yacoubi },
  year={ 2025 },
  eprint={ 2511.05841 },
  archivePrefix={arXiv},
  primaryClass={ cs.CV },
  url={ https://arxiv.org/abs/2511.05841 }
}


\bibitem{arxiv251105934}
@misc{das2025ad-dae:,
  title={ AD-DAE: Unsupervised Modeling of Longitudinal Alzheimer's Disease Progression with Diffusion Auto-Encoder },
  author={ Ayantika Das and Arunima Sarkar and Keerthi Ram and Mohanasankar Sivaprakasam },
  year={ 2025 },
  eprint={ 2511.05934 },
  archivePrefix={arXiv},
  primaryClass={ cs.CV },
  url={ https://arxiv.org/abs/2511.05934 }
}


\bibitem{arxiv251106215}
@misc{su2025explicit,
  title={ Explicit Knowledge-Guided In-Context Learning for Early Detection of Alzheimer's Disease },
  author={ Puzhen Su and Yongzhu Miao and Chunxi Guo and Jintao Tang and Shasha Li and Ting Wang },
  year={ 2025 },
  eprint={ 2511.06215 },
  archivePrefix={arXiv},
  primaryClass={ cs.CL },
  url={ https://arxiv.org/abs/2511.06215 }
}


\bibitem{arxiv251106681}
@misc{hou2025an,
  title={ An Adaptive Machine Learning Triage Framework for Predicting Alzheimer's Disease Progression },
  author={ Richard Hou and Shengpu Tang and Wei Jin },
  year={ 2025 },
  eprint={ 2511.06681 },
  archivePrefix={arXiv},
  primaryClass={ cs.LG },
  url={ https://arxiv.org/abs/2511.06681 }
}


\bibitem{arxiv251106826}
@misc{su2025beyond,
  title={ Beyond Plain Demos: A Demo-centric Anchoring Paradigm for In-Context Learning in Alzheimer's Disease Detection },
  author={ Puzhen Su and Haoran Yin and Yongzhu Miao and Jintao Tang and Shasha Li and Ting Wang },
  year={ 2025 },
  eprint={ 2511.06826 },
  archivePrefix={arXiv},
  primaryClass={ cs.CL },
  url={ https://arxiv.org/abs/2511.06826 }
}


\bibitem{arxiv251107313}
@misc{jacobson2025de-individualizing,
  title={ De-Individualizing fMRI Signals via Mahalanobis Whitening and Bures Geometry },
  author={ Aaron Jacobson and Tingting Dan and Martin Styner and Guorong Wu and Shahar Kovalsky and Caroline Moosmueller },
  year={ 2025 },
  eprint={ 2511.07313 },
  archivePrefix={arXiv},
  primaryClass={ q-bio.NC },
  url={ https://arxiv.org/abs/2511.07313 }
}


\bibitem{arxiv251108132}
@misc{zolnoori2025national,
  title={ National Institute on Aging PREPARE Challenge: Early Detection of Cognitive Impairment Using Speech -- The SpeechCARE Solution },
  author={ Maryam Zolnoori and Hossein Azadmaleki and Yasaman Haghbin and Ali Zolnour and Mohammad Javad Momeni Nezhad and Sina Rashidi and Mehdi Naserian and Elyas Esmaeili and Sepehr Karimi Arpanahi },
  year={ 2025 },
  eprint={ 2511.08132 },
  archivePrefix={arXiv},
  primaryClass={ cs.AI },
  url={ https://arxiv.org/abs/2511.08132 }
}


\bibitem{arxiv251108663}
@misc{ahmed20253d-tda,
  title={ 3D-TDA -- Topological feature extraction from 3D images for Alzheimer's disease classification },
  author={ Faisal Ahmed and Taymaz Akan and Fatih Gelir and Owen T. Carmichael and Elizabeth A. Disbrow and Steven A. Conrad and Mohammad A. N. Bhuiyan },
  year={ 2025 },
  eprint={ 2511.08663 },
  archivePrefix={arXiv},
  primaryClass={ eess.IV },
  url={ https://arxiv.org/abs/2511.08663 }
}


\bibitem{arxiv251108847}
@misc{li2025data-driven,
  title={ Data-driven spatiotemporal modeling reveals personalized trajectories of cortical atrophy in Alzheimer's disease },
  author={ Chunyan Li and Yutong Mao and Xiao Liu and Wenrui Hao },
  year={ 2025 },
  eprint={ 2511.08847 },
  archivePrefix={arXiv},
  primaryClass={ q-bio.NC },
  url={ https://arxiv.org/abs/2511.08847 }
}


\bibitem{arxiv251109949}
@misc{he2025imaging,
  title={ Imaging the Topology of Dynamic Brain Connectivity },
  author={ Peilin He and Tananun Songdechakraiwut },
  year={ 2025 },
  eprint={ 2511.09949 },
  archivePrefix={arXiv},
  primaryClass={ q-bio.NC },
  url={ https://arxiv.org/abs/2511.09949 }
}


\bibitem{arxiv251110890}
@misc{he2025llm,
  title={ LLM enhanced graph inference for long-term disease progression modelling },
  author={ Tiantian He and An Zhao and Elinor Thompson and Anna Schroder and Ahmed Abdulaal and Frederik Barkhof and Daniel C. Alexander },
  year={ 2025 },
  eprint={ 2511.10890 },
  archivePrefix={arXiv},
  primaryClass={ cs.AI },
  url={ https://arxiv.org/abs/2511.10890 }
}


\bibitem{arxiv251111996}
@misc{wu2025graphical,
  title={ Graphical Model-based Inference on Persistent Homology },
  author={ Zitian Wu and Arkaprava Roy and Leo L. Duan },
  year={ 2025 },
  eprint={ 2511.11996 },
  archivePrefix={arXiv},
  primaryClass={ stat.ME },
  url={ https://arxiv.org/abs/2511.11996 }
}


\bibitem{arxiv251112451}
@misc{bashit2025a,
  title={ A Multicollinearity-Aware Signal-Processing Framework for Cross-\$β\$ Identification via X-ray Scattering of Alzheimer's Tissue },
  author={ Abdullah Al Bashit and Prakash Nepal and Lee Makowski },
  year={ 2025 },
  eprint={ 2511.12451 },
  archivePrefix={arXiv},
  primaryClass={ eess.IV },
  url={ https://arxiv.org/abs/2511.12451 }
}


\bibitem{arxiv251113911}
@misc{tassopoulou2025uncertainty-calibrated,
  title={ Uncertainty-Calibrated Prediction of Randomly-Timed Biomarker Trajectories with Conformal Bands },
  author={ Vasiliki Tassopoulou and Charis Stamouli and Haochang Shou and George J. Pappas and Christos Davatzikos },
  year={ 2025 },
  eprint={ 2511.13911 },
  archivePrefix={arXiv},
  primaryClass={ stat.ML },
  url={ https://arxiv.org/abs/2511.13911 }
}


\bibitem{arxiv251114417}
@misc{redondo2025nonlinear,
  title={ Nonlinear Coherence for Vector Time Series: Defining Region-to-Region Functional Brain Connectivity },
  author={ Paolo Victor Redondo and Raphaël Huser and Hernando Ombao },
  year={ 2025 },
  eprint={ 2511.14417 },
  archivePrefix={arXiv},
  primaryClass={ stat.AP },
  url={ https://arxiv.org/abs/2511.14417 }
}


\bibitem{arxiv251114588}
@misc{machnio2025deep,
  title={ Deep Learning-Based Regional White Matter Hyperintensity Mapping as a Robust Biomarker for Alzheimer's Disease },
  author={ Julia Machnio and Mads Nielsen and Mostafa Mehdipour Ghazi },
  year={ 2025 },
  eprint={ 2511.14588 },
  archivePrefix={arXiv},
  primaryClass={ cs.CV },
  url={ https://arxiv.org/abs/2511.14588 }
}


\bibitem{arxiv251114601}
@misc{putera2025mri,
  title={ MRI Embeddings Complement Clinical Predictors for Cognitive Decline Modeling in Alzheimer's Disease Cohorts },
  author={ Nathaniel Putera and Daniel Vilet Rodríguez and Noah Videcrantz and Julia Machnio and Mostafa Mehdipour Ghazi },
  year={ 2025 },
  eprint={ 2511.14601 },
  archivePrefix={arXiv},
  primaryClass={ cs.CV },
  url={ https://arxiv.org/abs/2511.14601 }
}


\bibitem{arxiv251114922}
@misc{peddi2025integrating,
  title={ Integrating Causal Inference with Graph Neural Networks for Alzheimer's Disease Analysis },
  author={ Pranay Kumar Peddi and Dhrubajyoti Ghosh },
  year={ 2025 },
  eprint={ 2511.14922 },
  archivePrefix={arXiv},
  primaryClass={ cs.LG },
  url={ https://arxiv.org/abs/2511.14922 }
}


\bibitem{arxiv251115151}
@misc{zhou2025dcl-se:,
  title={ DCL-SE: Dynamic Curriculum Learning for Spatiotemporal Encoding of Brain Imaging },
  author={ Meihua Zhou and Xinyu Tong and Jiarui Zhao and Min Cheng and Li Yang and Lei Tian and Nan Wan },
  year={ 2025 },
  eprint={ 2511.15151 },
  archivePrefix={arXiv},
  primaryClass={ cs.CV },
  url={ https://arxiv.org/abs/2511.15151 }
}


\bibitem{arxiv251115188}
@misc{jalal2025brainrotvit:,
  title={ BrainRotViT: Transformer-ResNet Hybrid for Explainable Modeling of Brain Aging from 3D sMRI },
  author={ Wasif Jalal and Md Nafiu Rahman and Atif Hasan Rahman and M. Sohel Rahman },
  year={ 2025 },
  eprint={ 2511.15188 },
  archivePrefix={arXiv},
  primaryClass={ cs.CV },
  url={ https://arxiv.org/abs/2511.15188 }
}


\bibitem{arxiv251118238}
@misc{samadzelkava2025detecting,
  title={ Detecting Discontinuities in the Topology of Alzheimers gene Co-expression },
  author={ Aya Samadzelkava },
  year={ 2025 },
  eprint={ 2511.18238 },
  archivePrefix={arXiv},
  primaryClass={ q-bio.QM },
  url={ https://arxiv.org/abs/2511.18238 }
}


\bibitem{arxiv251120154}
@misc{hong2025alzheimers,
  title={ Alzheimers Disease Progression Prediction Based on Manifold Mapping of Irregularly Sampled Longitudinal Data },
  author={ Xin Hong and Ying Shi and Yinhao Li and Yen-Wei Chen },
  year={ 2025 },
  eprint={ 2511.20154 },
  archivePrefix={arXiv},
  primaryClass={ cs.CV },
  url={ https://arxiv.org/abs/2511.20154 }
}


\bibitem{arxiv251120704}
@misc{moslemi2025pretraining,
  title={ Pretraining Transformer-Based Models on Diffusion-Generated Synthetic Graphs for Alzheimer's Disease Prediction },
  author={ Abolfazl Moslemi and Hossein Peyvandi },
  year={ 2025 },
  eprint={ 2511.20704 },
  archivePrefix={arXiv},
  primaryClass={ cs.LG },
  url={ https://arxiv.org/abs/2511.20704 }
}


\bibitem{arxiv251121057}
@misc{hong2025long-term,
  title={ Long-Term Alzheimers Disease Prediction: A Novel Image Generation Method Using Temporal Parameter Estimation with Normal Inverse Gamma Distribution on Uneven Time Series },
  author={ Xin Hong and Xinze Sun and Yinhao Li and Yen-Wei Chen },
  year={ 2025 },
  eprint={ 2511.21057 },
  archivePrefix={arXiv},
  primaryClass={ cs.CV },
  url={ https://arxiv.org/abs/2511.21057 }
}


\bibitem{arxiv251121114}
@misc{honga2025deformation-aware,
  title={ Deformation-aware Temporal Generation for Early Prediction of Alzheimers Disease },
  author={ Xin Honga and Jie Lin and Minghui Wang },
  year={ 2025 },
  eprint={ 2511.21114 },
  archivePrefix={arXiv},
  primaryClass={ cs.CV },
  url={ https://arxiv.org/abs/2511.21114 }
}


\bibitem{arxiv251121530}
@misc{hong2025the,
  title={ The Age-specific Alzheimer 's Disease Prediction with Characteristic Constraints in Nonuniform Time Span },
  author={ Xin Hong and Kaifeng Huang },
  year={ 2025 },
  eprint={ 2511.21530 },
  archivePrefix={arXiv},
  primaryClass={ cs.CV },
  url={ https://arxiv.org/abs/2511.21530 }
}


\bibitem{arxiv251121724}
@misc{sun2025ad-cdo:,
  title={ AD-CDO: A Lightweight Ontology for Representing Eligibility Criteria in Alzheimer's Disease Clinical Trials },
  author={ Zenan Sun and Rashmie Abeysinghe and Xiaojin Li and Xinyue Hu and Licong Cui and Guo-Qiang Zhang and Jiang Bian and Cui Tao },
  year={ 2025 },
  eprint={ 2511.21724 },
  archivePrefix={arXiv},
  primaryClass={ cs.CL },
  url={ https://arxiv.org/abs/2511.21724 }
}


\bibitem{arxiv251122102}
@misc{crête2025mri-based,
  title={ MRI-Based Brain Age Estimation with Supervised Contrastive Learning of Continuous Representation },
  author={ Simon Joseph Clément Crête and Marta Kersten-Oertel and Yiming Xiao },
  year={ 2025 },
  eprint={ 2511.22102 },
  archivePrefix={arXiv},
  primaryClass={ cs.CV },
  url={ https://arxiv.org/abs/2511.22102 }
}


\bibitem{arxiv251122774}
@misc{khatooni2025alzheimer's,
  title={ Alzheimer's Disease Prediction Using EffNetViTLoRA and BiLSTM with Multimodal Longitudinal MRI Data },
  author={ Mahdieh Behjat Khatooni and Mohsen Soryani },
  year={ 2025 },
  eprint={ 2511.22774 },
  archivePrefix={arXiv},
  primaryClass={ cs.CV },
  url={ https://arxiv.org/abs/2511.22774 }
}


\bibitem{arxiv251200269}
@misc{wang2025usb:,
  title={ USB: Unified Synthetic Brain Framework for Bidirectional Pathology-Healthy Generation and Editing },
  author={ Jun Wang and Peirong Liu },
  year={ 2025 },
  eprint={ 2512.00269 },
  archivePrefix={arXiv},
  primaryClass={ cs.CV },
  url={ https://arxiv.org/abs/2512.00269 }
}


\bibitem{arxiv251201367}
@misc{jiang2025a,
  title={ A Fine Evaluation Method for Cube Copying Test for Early Detection of Alzheimer's Disease },
  author={ Xinyu Jiang and Cuiyun Gao and Wenda Huang and Yiyang Jiang and Binwen Luo and Yuxin Jiang and Mengting Wang and Haoran Wen and Yang Zhao and Xuemei Chen and Songqun Huang },
  year={ 2025 },
  eprint={ 2512.01367 },
  archivePrefix={arXiv},
  primaryClass={ cs.LG },
  url={ https://arxiv.org/abs/2512.01367 }
}


\bibitem{arxiv251203467}
@misc{hao2025bayesian,
  title={ Bayesian Event-Based Model for Disease Subtype and Stage Inference },
  author={ Hongtao Hao and Joseph L. Austerweil },
  year={ 2025 },
  eprint={ 2512.03467 },
  archivePrefix={arXiv},
  primaryClass={ cs.LG },
  url={ https://arxiv.org/abs/2512.03467 }
}


\bibitem{arxiv251203912}
@misc{xu2025parsimonious,
  title={ Parsimonious Clustering of Covariance Matrices },
  author={ Yixi Xu and Yi Zhao },
  year={ 2025 },
  eprint={ 2512.03912 },
  archivePrefix={arXiv},
  primaryClass={ stat.ME },
  url={ https://arxiv.org/abs/2512.03912 }
}


\bibitem{arxiv251205814}
@misc{zhu2025ug-fedda:,
  title={ UG-FedDA: Uncertainty-Guided Federated Domain Adaptation for Multi-Center Alzheimer's Disease Detection },
  author={ Fubao Zhu and Zhanyuan Jia and Zhiguo Wang and Huan Huang and Danyang Sun and Chuang Han and Yanting Li and Jiaofen Nan and Chen Zhao and Weihua Zhou },
  year={ 2025 },
  eprint={ 2512.05814 },
  archivePrefix={arXiv},
  primaryClass={ cs.CV },
  url={ https://arxiv.org/abs/2512.05814 }
}


\bibitem{arxiv251205993}
@misc{verma2025domain-specific,
  title={ Domain-Specific Foundation Model Improves AI-Based Analysis of Neuropathology },
  author={ Ruchika Verma and Shrishtee Kandoi and Robina Afzal and Shengjia Chen and Jannes Jegminat and Michael W. Karlovich and Melissa Umphlett and Timothy E. Richardson and Kevin Clare and Quazi Hossain and Jorge Samanamud and Phyllis L. Faust and Elan D. Louis and Ann C. McKee and Thor D. Stein and Jonathan D. Cherry and Jesse Mez and Anya C. McGoldrick and Dalilah D. Quintana Mora and Melissa J. Nirenberg and Ruth H. Walker and Yolfrankcis Mendez and Susan Morgello and Dennis W. Dickson and Melissa E. Murray and Carlos Cordon-Cardo and Nadejda M. Tsankova and Jamie M. Walker and Diana K. Dangoor and Stephanie McQuillan and Emma L. Thorn and Claudia De Sanctis and Shuying Li and Thomas J. Fuchs and Kurt Farrell and John F. Crary and Gabriele Campanella },
  year={ 2025 },
  eprint={ 2512.05993 },
  archivePrefix={arXiv},
  primaryClass={ cs.CV },
  url={ https://arxiv.org/abs/2512.05993 }
}


\bibitem{arxiv251206134}
@misc{hrusanov2025physics-informed,
  title={ Physics-Informed Neural Koopman Machine for Interpretable Longitudinal Personalized Alzheimer's Disease Forecasting },
  author={ Georgi Hrusanov and Duy-Thanh Vu and Duy-Cat Can and Sophie Tascedda and Margaret Ryan and Julien Bodelet and Katarzyna Koscielska and Carsten Magnus and Oliver Y. Chén },
  year={ 2025 },
  eprint={ 2512.06134 },
  archivePrefix={arXiv},
  primaryClass={ cs.LG },
  url={ https://arxiv.org/abs/2512.06134 }
}


\bibitem{arxiv251207061}
@misc{maity2025alterations,
  title={ Alterations of brain tissue structural complexity and disorder in Alzheimer's disease (AD): Fractal, multifractal, fractal transformation, and disorder strength analyses },
  author={ Santanu Maity and Mousa Alrubayan and Mohammad Moshahid Khan and Prabhakar Pradhan },
  year={ 2025 },
  eprint={ 2512.07061 },
  archivePrefix={arXiv},
  primaryClass={ physics.med-ph },
  url={ https://arxiv.org/abs/2512.07061 }
}


\bibitem{arxiv251209217}
@misc{namadi2025access,
  title={ Access to healthcare for people with Alzheimer's Diseases and related dementias },
  author={ Saeed Saleh Namadi and Jie Chen and Deb Niemeier },
  year={ 2025 },
  eprint={ 2512.09217 },
  archivePrefix={arXiv},
  primaryClass={ stat.AP },
  url={ https://arxiv.org/abs/2512.09217 }
}


\bibitem{arxiv251210966}
@misc{zhuang2025multimodal,
  title={ Multimodal Fusion of Regional Brain Experts for Interpretable Alzheimer's Disease Diagnosis },
  author={ Farica Zhuang and Dinara Aliyeva and Shu Yang and Zixuan Wen and Duy Duong-Tran and Christos Davatzikos and Tianlong Chen and Song Wang and Li Shen },
  year={ 2025 },
  eprint={ 2512.10966 },
  archivePrefix={arXiv},
  primaryClass={ cs.LG },
  url={ https://arxiv.org/abs/2512.10966 }
}


\bibitem{arxiv251213346}
@misc{lu2025beyond,
  title={ Beyond Missing Data: Questionnaire Uncertainty Responses as Early Digital Biomarkers of Cognitive Decline and Neurodegenerative Diseases },
  author={ Yukun Lu and Bingjie Li and Zhigang Yao },
  year={ 2025 },
  eprint={ 2512.13346 },
  archivePrefix={arXiv},
  primaryClass={ stat.AP },
  url={ https://arxiv.org/abs/2512.13346 }
}


\bibitem{arxiv251213685}
@misc{phelps2025beyond,
  title={ Beyond surface form: A pipeline for semantic analysis in Alzheimer's Disease detection from spontaneous speech },
  author={ Dylan Phelps and Rodrigo Wilkens and Edward Gow-Smith and Lilian Hubner and Bárbara Malcorra and César Rennó-Costa and Marco Idiart and Maria-Cruz Villa-Uriol and Aline Villavicencio },
  year={ 2025 },
  eprint={ 2512.13685 },
  archivePrefix={arXiv},
  primaryClass={ cs.CL },
  url={ https://arxiv.org/abs/2512.13685 }
}


\bibitem{arxiv251213724}
@misc{noori2025graph,
  title={ Graph AI generates neurological hypotheses validated in molecular, organoid, and clinical systems },
  author={ Ayush Noori and Joaquín Polonuer and Katharina Meyer and Bogdan Budnik and Shad Morton and Xinyuan Wang and Sumaiya Nazeen and Yingnan He and Iñaki Arango and Lucas Vittor and Matthew Woodworth and Richard C. Krolewski and Michelle M. Li and Ninning Liu and Tushar Kamath and Evan Macosko and Dylan Ritter and Jalwa Afroz and Alexander B. H. Henderson and Lorenz Studer and Samuel G. Rodriques and Andrew White and Noa Dagan and David A. Clifton and George M. Church and Sudeshna Das and Jenny M. Tam and Vikram Khurana and Marinka Zitnik },
  year={ 2025 },
  eprint={ 2512.13724 },
  archivePrefix={arXiv},
  primaryClass={ q-bio.QM },
  url={ https://arxiv.org/abs/2512.13724 }
}


\bibitem{arxiv251213747}
@misc{dai2025why,
  title={ Why Text Prevails: Vision May Undermine Multimodal Medical Decision Making },
  author={ Siyuan Dai and Lunxiao Li and Kun Zhao and Eardi Lila and Paul K. Crane and Heng Huang and Dongkuan Xu and Haoteng Tang and Liang Zhan },
  year={ 2025 },
  eprint={ 2512.13747 },
  archivePrefix={arXiv},
  primaryClass={ cs.CV },
  url={ https://arxiv.org/abs/2512.13747 }
}


\bibitem{arxiv251214026}
@misc{fu2025unleashing,
  title={ Unleashing the Power of Image-Tabular Self-Supervised Learning via Breaking Cross-Tabular Barriers },
  author={ Yibing Fu and Yunpeng Zhao and Zhitao Zeng and Cheng Chen and Yueming Jin },
  year={ 2025 },
  eprint={ 2512.14026 },
  archivePrefix={arXiv},
  primaryClass={ cs.CV },
  url={ https://arxiv.org/abs/2512.14026 }
}


\bibitem{arxiv251214491}
@misc{lu2025sparse,
  title={ Sparse Multi-Modal Transformer with Masking for Alzheimer's Disease Classification },
  author={ Cheng-Han Lu and Pei-Hsuan Tsai },
  year={ 2025 },
  eprint={ 2512.14491 },
  archivePrefix={arXiv},
  primaryClass={ cs.AI },
  url={ https://arxiv.org/abs/2512.14491 }
}


\bibitem{arxiv251215246}
@misc{vinh2025enhancing,
  title={ Enhancing Alzheimer's Detection through Late Fusion of Multi-Modal EEG Features },
  author={ Nguyen Thanh Vinh and Manoj Vishwanath and Thinh Nguyen-Quang and Nguyen Viet Ha and Bui Thanh Tung and Huy-Dung Han and Nguyen Quang Linh and Nguyen Hai Linh and Hung Cao },
  year={ 2025 },
  eprint={ 2512.15246 },
  archivePrefix={arXiv},
  primaryClass={ eess.SP },
  url={ https://arxiv.org/abs/2512.15246 }
}


\bibitem{arxiv251215947}
@misc{kim2025mcr-vqgan:,
  title={ MCR-VQGAN: A Scalable and Cost-Effective Tau PET Synthesis Approach for Alzheimer's Disease Imaging },
  author={ Jin Young Kim and Jeremy Hudson and Jeongchul Kim and Qing Lyu and Christopher T. Whitlow },
  year={ 2025 },
  eprint={ 2512.15947 },
  archivePrefix={arXiv},
  primaryClass={ eess.IV },
  url={ https://arxiv.org/abs/2512.15947 }
}


\bibitem{arxiv251216184}
@misc{yang2025a,
  title={ A Multimodal Approach to Alzheimer's Diagnosis: Geometric Insights from Cube Copying and Cognitive Assessments },
  author={ Jaeho Yang and Kijung Yoon },
  year={ 2025 },
  eprint={ 2512.16184 },
  archivePrefix={arXiv},
  primaryClass={ cs.LG },
  url={ https://arxiv.org/abs/2512.16184 }
}


\bibitem{arxiv251216964}
@misc{ahmed2025colormap-enhanced,
  title={ Colormap-Enhanced Vision Transformers for MRI-Based Multiclass (4-Class) Alzheimer's Disease Classification },
  author={ Faisal Ahmed },
  year={ 2025 },
  eprint={ 2512.16964 },
  archivePrefix={arXiv},
  primaryClass={ eess.IV },
  url={ https://arxiv.org/abs/2512.16964 }
}


\bibitem{arxiv251217276}
@misc{moayedikia2025alzheimer's,
  title={ Alzheimer's Disease Brain Network Mining },
  author={ Alireza Moayedikia and Sara Fin },
  year={ 2025 },
  eprint={ 2512.17276 },
  archivePrefix={arXiv},
  primaryClass={ cs.LG },
  url={ https://arxiv.org/abs/2512.17276 }
}


\bibitem{arxiv251218403}
@misc{li2025bayesian,
  title={ Bayesian Brain Edge-Based Connectivity (BBeC): a Bayesian model for brain edge-based connectivity inference },
  author={ Zijing Li and Chenhao Zeng and Shufei Ge },
  year={ 2025 },
  eprint={ 2512.18403 },
  archivePrefix={arXiv},
  primaryClass={ stat.ME },
  url={ https://arxiv.org/abs/2512.18403 }
}


\bibitem{arxiv251218986}
@misc{zhao2025r-genima:,
  title={ R-GenIMA: Integrating Neuroimaging and Genetics with Interpretable Multimodal AI for Alzheimer's Disease Progression },
  author={ Kun Zhao and Siyuan Dai and Yingying Zhang and Guodong Liu and Pengfei Gu and Chenghua Lin and Paul M. Thompson and Alex Leow and Heng Huang and Lifang He and Liang Zhan and Haoteng Tang },
  year={ 2025 },
  eprint={ 2512.18986 },
  archivePrefix={arXiv},
  primaryClass={ cs.LG },
  url={ https://arxiv.org/abs/2512.18986 }
}


\bibitem{arxiv251219099}
@misc{moayedikia2025dual,
  title={ Dual Model Deep Learning for Alzheimer Prognostication },
  author={ Alireza Moayedikia and Sara Fin and Uffe Kock Wiil },
  year={ 2025 },
  eprint={ 2512.19099 },
  archivePrefix={arXiv},
  primaryClass={ cs.LG },
  url={ https://arxiv.org/abs/2512.19099 }
}


\bibitem{arxiv251220948}
@misc{dong2025foundation,
  title={ Foundation Model-based Evaluation of Neuropsychiatric Disorders: A Lifespan-Inclusive, Multi-Modal, and Multi-Lingual Study },
  author={ Zhongren Dong and Haotian Guo and Weixiang Xu and Huan Zhao and Zixing Zhang },
  year={ 2025 },
  eprint={ 2512.20948 },
  archivePrefix={arXiv},
  primaryClass={ cs.CL },
  url={ https://arxiv.org/abs/2512.20948 }
}


\bibitem{arxiv251223093}
@misc{drishti2025cogniscope:,
  title={ Cogniscope: Modeling Social Media Interactions as Digital Biomarkers for Early Detection of Cognitive Decline },
  author={ Ananya Drishti and Mahfuza Farooque },
  year={ 2025 },
  eprint={ 2512.23093 },
  archivePrefix={arXiv},
  primaryClass={ cs.HC },
  url={ https://arxiv.org/abs/2512.23093 }
}


\bibitem{arxiv251223441}
@misc{emre2025stochastic,
  title={ Stochastic Siamese MAE Pretraining for Longitudinal Medical Images },
  author={ Taha Emre and Arunava Chakravarty and Thomas Pinetz and Dmitrii Lachinov and Martin J. Menten and Hendrik Scholl and Sobha Sivaprasad and Daniel Rueckert and Andrew Lotery and Stefan Sacu and Ursula Schmidt-Erfurth and Hrvoje Bogunović },
  year={ 2025 },
  eprint={ 2512.23441 },
  archivePrefix={arXiv},
  primaryClass={ cs.LG },
  url={ https://arxiv.org/abs/2512.23441 }
}


\bibitem{arxiv260100143}
@misc{qu2026methconvtransformer:,
  title={ MethConvTransformer: A Deep Learning Framework for Cross-Tissue Alzheimer's Disease Detection },
  author={ Gang Qu and Guanghao Li and Zhongming Zhao },
  year={ 2026 },
  eprint={ 2601.00143 },
  archivePrefix={arXiv},
  primaryClass={ q-bio.GN },
  url={ https://arxiv.org/abs/2601.00143 }
}


\end{thebibliography}
\end{document}